\chapter{Kết luận và hướng phát triển}
\label{chap:conclusion}

Chương này tổng kết lại những kết quả đạt được của dự án nghiên cứu và phát triển hệ thống Self-Evolving AI Agents trên nền tảng NIKA, đồng thời thảo luận thẳng thắn về những hạn chế hiện tại và đề xuất định hướng phát triển trong tương lai.

\section{Kết quả đạt được của dự án}
\label{sec:accomplishments}

Sau quá trình nghiên cứu và thực nghiệm nghiêm túc, dự án đã đạt được một số kết quả quan trọng sau:
\begin{itemize}
    \item \textbf{Nghiên cứu lý thuyết có trọng tâm:} Tổng hợp và hệ thống hóa hai hướng nghiên cứu liên quan trực tiếp đến đề tài: procedural memory trong Skill-Pro và tinh chỉnh tài liệu công cụ trong DRAFT.
    \item \textbf{Xây dựng kiến trúc hệ thống hoàn chỉnh:} Tích hợp thành công AI Agent với môi trường giả lập mạng Kathará thông qua giao thức Model Context Protocol (MCP) chuẩn hóa. Hệ thống vận hành tốt trên 3 luồng chẩn đoán phổ biến (ReAct, Plan \& Execute, Reflexion).
    \item \textbf{Hiện thực hóa thành công các mô-đun tiến hóa:}
        \begin{enumerate}
            \item Mô-đun Procedural Memory lưu trữ \texttt{ProceduralSkill} trong file JSON (\texttt{runtime/memory/\{bank\_id\}/skills.json}), hỗ trợ truy xuất kỹ năng theo ngữ cảnh và tiến hóa kỹ năng ngoại tuyến.
            \item Mô-đun Tool Refinement tự tinh chỉnh mô tả các công cụ MCP nguyên thủy thông qua vòng lặp Explorer--Analyzer--Rewriter. Trạng thái được lưu tại thư mục \texttt{runtime/tool\_evolution/} với cơ chế đóng băng adaptive khi hội tụ.
        \end{enumerate}
    \item \textbf{Đánh giá thực nghiệm chi tiết:} Kiểm chứng hệ thống qua hai lần chạy benchmark trên 125 case chung của NIKA. Kết quả cho thấy Evolved Agent đạt incident success rate 28.0\% (35/125), tăng 34.6\% so với Baseline Agent 20.8\% (26/125). RCA F1 cải thiện từ 0.304 lên 0.387, số bước chẩn đoán giảm 16.7\%, số lần gọi công cụ giảm 18.5\% và tool error rate giảm từ 0.65\% xuống 0.00\%.
\end{itemize}

\section{Những hạn chế hiện tại}
\label{sec:limitations}

Mặc dù đạt được những kết quả khả quan, hệ thống vẫn tồn tại một số hạn chế cần khắc phục:
\begin{itemize}
    \item \textbf{Khả năng tổng quát hóa giữa các kịch bản còn hạn chế:} Các kỹ năng và tài liệu công cụ được tiến hóa từ những episode trước có thể thiên về cấu trúc topology, tên thiết bị hoặc giao thức đã gặp. Vì vậy, khi chuyển sang các kịch bản phức tạp hơn như BGP/OSPF, SDN, P4 hoặc các lỗi resource contention, agent vẫn có thể chọn sai hướng điều tra, gọi công cụ với tham số chưa phù hợp hoặc chưa suy luận đủ sâu từ bằng chứng thu thập được.
    \item \textbf{Tăng chi phí Token và thời gian chạy:} Việc tải thêm tài liệu công cụ đã được tinh chỉnh và các kỹ năng từ Procedural Memory vào ngữ cảnh prompt làm tổng token tăng 12.4\% (từ 16.641M lên 18.699M token), chủ yếu đến từ input tokens. Tổng thời gian chạy tăng 7.2\% (từ 207.7 lên 222.6 phút).
    \item \textbf{Chất lượng tiến hóa phụ thuộc vào LLM:} Cả Semantic Gradient trong Procedural Memory lẫn Rewriter trong Tool Refinement đều dựa vào LLM để sinh đề xuất cải tiến. Khi LLM sinh gradient hoặc tài liệu sai lệch, hệ thống phải dùng fallback xác định, làm giảm chất lượng học hỏi.
\end{itemize}

\section{Định hướng phát triển tương lai}
\label{sec:future_work}

Từ những kết quả và hạn chế trên, dự án định hướng một số nội dung nghiên cứu tiếp theo bao gồm:
\begin{itemize}
    \item \textbf{Cải thiện khả năng tổng quát hóa (Generalization):} Nghiên cứu các cơ chế trừu tượng hóa ngữ nghĩa mạnh mẽ hơn trong Procedural Memory và Tool Refinement, giúp kỹ năng chẩn đoán và tài liệu công cụ tinh chỉnh dễ dàng thích ứng với topology, tên thiết bị, giao thức và nhóm lỗi mới chưa gặp.
    \item \textbf{Tích hợp cơ chế Human-in-the-loop:} Bổ sung giao diện cho phép kỹ sư mạng tham gia phản hồi, đánh giá và phê duyệt các kỹ năng mới được PPO Gate chấp nhận cũng như tài liệu công cụ đã tinh chỉnh trước khi ghi nhận chính thức vào bộ nhớ.
    \item \textbf{Tối ưu hóa hiệu năng và chi phí:} Tinh chỉnh các mô hình ngôn ngữ lớn chuyên biệt cho miền mạng (Domain-specific LLMs) để giảm kích thước prompt chẩn đoán, tối ưu hóa thời gian thực thi các Kathará labs, và giảm chi phí token khi tải kỹ năng từ Procedural Memory.
    \item \textbf{Mở rộng kịch bản sự cố:} Đánh giá agent tự tiến hóa trên các kịch bản mạng có quy mô lớn hơn (topology tier m, l) và các sự cố phức tạp hơn (multi-fault incidents với nhiều nguyên nhân đồng thời).
\end{itemize}
